\section{Conclusão}
\label{chap:conclusao}
O estudo demonstrou que, para o dataset avaliado, um modelo linear simples (MQ), quando combinado com uma estratégia de pré-processamento robusta, pode superar o desempenho de modelos não-lineares mais complexos como as redes MLP. 

Este resultado desafia a premissa de que a complexidade crescente do modelo invariavelmente leva a um melhor desempenho e destaca a importância crítica da preparação dos dados. A análise de diagnóstico revelou a presença de pontos influentes e heterocedasticidade, enquanto a não-normalidade universal dos resíduos aponta para características intrínsecas dos dados. 